% !TEX program = xelatex
\documentclass[11pt,a4paper]{article}

\usepackage{fontspec}
\usepackage{polyglossia}
\setmainlanguage{russian}
\setotherlanguage{english}
\setmainfont{DejaVu Serif}[Scale=0.9]
\setsansfont{DejaVu Sans}[Scale=0.9]
\setmonofont{DejaVu Sans Mono}[Scale=0.82]

\usepackage[margin=2.3cm]{geometry}
\usepackage{xcolor}
\usepackage{booktabs,longtable,tabularx,array,enumitem,pdflscape,microtype,ragged2e}
\usepackage{fancyhdr,titlesec}
\usepackage{tocloft}
\setlength{\cftsecnumwidth}{2em}
\setlength{\cftsubsecnumwidth}{3em}
\usepackage{listings}
\usepackage[most]{tcolorbox}
\usepackage[hidelinks]{hyperref}

\definecolor{accent}{HTML}{1F4E79}
\definecolor{okgreen}{HTML}{2E7D32}
\definecolor{warnred}{HTML}{B71C1C}
\hypersetup{colorlinks=true,linkcolor=accent,urlcolor=accent,citecolor=accent}
\Urlmuskip=0mu plus 1mu\relax
\makeatletter\g@addto@macro\UrlBreaks{\do\-\do\_\do\=\do\&}\makeatother
\setlength{\emergencystretch}{3em}

\setlist{itemsep=2pt,topsep=4pt,parsep=0pt}
\setlength{\parindent}{0pt}
\setlength{\parskip}{5pt}

\titleformat{\section}{\Large\sffamily\bfseries\color{accent}}{\thesection}{0.6em}{}
\titleformat{\subsection}{\large\sffamily\bfseries\color{accent}}{\thesubsection}{0.6em}{}
\titleformat{\subsubsection}{\normalsize\sffamily\bfseries}{\thesubsubsection}{0.6em}{}

\pagestyle{fancy}
\fancyhf{}
\fancyhead[L]{\small\sffamily Поиск научных статей: программный доступ}
\fancyhead[R]{\small\sffamily\thepage}
\renewcommand{\headrulewidth}{0.3pt}

\newcommand{\code}[1]{\texttt{\detokenize{#1}}}

% ---------- listings ----------
\lstdefinestyle{py}{
  language=Python,
  basicstyle=\ttfamily\footnotesize,
  keywordstyle=\color{accent}\bfseries,
  commentstyle=\color{gray}\itshape,
  stringstyle=\color{okgreen!70!black},
  breaklines=true,
  breakatwhitespace=false,
  frame=single,
  rulecolor=\color{black!25},
  backgroundcolor=\color{black!3},
  columns=fullflexible,
  keepspaces=true,
  showstringspaces=false,
  xleftmargin=4pt,
  xrightmargin=4pt,
  aboveskip=6pt,
  belowskip=4pt
}
\lstdefinestyle{plain}{
  basicstyle=\ttfamily\footnotesize,
  breaklines=true,
  frame=single,
  rulecolor=\color{black!25},
  backgroundcolor=\color{black!3},
  columns=fullflexible,
  keepspaces=true,
  showstringspaces=false,
  xleftmargin=4pt,
  xrightmargin=4pt,
  aboveskip=6pt,
  belowskip=4pt
}
\lstset{style=py}

% ---------- cards ----------
\newtcolorbox{syscard}[1]{
  breakable,
  enhanced,
  colback=black!2,
  colframe=accent!70,
  coltitle=white,
  fonttitle=\sffamily\bfseries,
  title={#1},
  left=6pt,right=6pt,top=4pt,bottom=2pt,
  boxrule=0.6pt,
  arc=2pt
}
\newcommand{\row}[2]{\noindent{\sffamily\bfseries\small\color{accent}#1.}\ #2\par\vspace{3pt}}

\newtcolorbox{warnbox}[1][Важно]{
  breakable,colback=warnred!5,colframe=warnred!70!black,
  coltitle=white,fonttitle=\sffamily\bfseries,title={#1},
  left=6pt,right=6pt,top=3pt,bottom=3pt,arc=2pt,boxrule=0.6pt
}
\newtcolorbox{notebox}[1][Замечание]{
  breakable,colback=accent!5,colframe=accent!70,
  coltitle=white,fonttitle=\sffamily\bfseries,title={#1},
  left=6pt,right=6pt,top=3pt,bottom=3pt,arc=2pt,boxrule=0.6pt
}

\title{\sffamily\bfseries\color{accent}
Поиск научных статей в разных системах:\\[2pt]
программный доступ, особенности, ограничения\\и блокировка веб-агентов}
\author{}
\date{10 ноября 2026\,г.}

\begin{document}
\maketitle
\thispagestyle{fancy}

\begin{abstract}
\noindent Документ описывает более двадцати систем поиска научной литературы с точки зрения автоматизации: на что ориентирована каждая система, как получить доступ из кода (API, OAI-PMH, дампы), какие нужны ключи, каков синтаксис запросов, какие действуют лимиты и как система относится к автоматическим клиентам, то есть блокирует ли она веб-агентов. Для основных систем приведены примеры на Python. В конце даны сводные таблицы, рекомендации по выбору источников, универсальный клиент с корректной обработкой лимитов и скрипт для самостоятельной проверки доступности систем из вашей сети.
\end{abstract}

\tableofcontents
\newpage

% =====================================================================
\section{Как читать этот документ}
% =====================================================================

\subsection{Методика и степень достоверности}

Сведения собраны 10 ноября 2026\,г. из официальной документации, страниц политик использования, анонсов изменений и вторичных источников (описания клиентских библиотек, форумы, обзоры). Лимиты и условия доступа у многих сервисов менялись в 2025--2026 годах, поэтому для каждой системы указан \emph{статус данных}:

\begin{itemize}
  \item \textbf{официальные} --- подтверждено документацией или анонсом самого сервиса;
  \item \textbf{вторичные} --- взято из сторонних описаний, официальную страницу подтвердить не удалось;
  \item \textbf{противоречивые} --- разные источники называют разные цифры; нужно проверить на странице сервиса или самим запросом.
\end{itemize}

\begin{warnbox}[Что не было проверено]
Прямых тестовых запросов к сайтам я не делал: из рабочей среды исходящая сеть для оболочки закрыта, а попытка получить \code{robots.txt} нескольких сайтов штатным загрузчиком была отклонена. Поэтому выводы о блокировке веб-агентов опираются на документацию, политики сервисов и описания реальных блокировок, но не на живой тест. В разделе~\ref{sec:probe} приведён скрипт, который проверит всё это из вашей сети за пару минут.
\end{warnbox}

\subsection{Основные термины}

\begin{description}[leftmargin=2.2cm,style=nextline]
  \item[API-ключ] Токен, который выдают после бесплатной или платной регистрации. Обычно повышает лимиты и позволяет сервису связаться с вами.
  \item[Polite pool] Приоритетный пул у Crossref и OpenAlex (ранее): запрос содержит контактный e-mail (\code{mailto=}), за это дают более высокие лимиты.
  \item[Пагинация курсором] Выдача страницами через токен (\code{cursor}, \code{cursorMark}, \code{scroll}); в отличие от \code{offset} не упирается в потолок глубины.
  \item[OAI-PMH] Стандартный протокол массового сбора метаданных из репозиториев.
  \item[Снапшот (дамп)] Полная копия базы для локальной работы; когда нужны миллионы записей, это правильнее, чем API.
  \item[Веб-агент / бот] Автоматический клиент. Часть сервисов считает ботом любой клиент без браузера и блокирует его на уровне HTML-страниц, но разрешает на официальном API.
\end{description}

\subsection{Общий принцип выбора канала доступа}

Для автоматизации всегда следует идти по лестнице от «разрешённого» к «терпимому», а не наоборот:

\begin{enumerate}
  \item \textbf{Официальный API} --- для поиска и точечных запросов.
  \item \textbf{Снапшот или OAI-PMH} --- для массовой выгрузки.
  \item \textbf{Разбор HTML} --- только если система явно не запрещает это и других путей нет. Для большинства систем из этого документа HTML-скрейпинг запрещён или приводит к блокировке.
\end{enumerate}

\begin{notebox}[Граница, которую документ не переходит]
Здесь не описаны способы обхода CAPTCHA, ротации IP-адресов, подмены User-Agent и прочие приёмы обхода защиты. Если сервис блокирует автоматический доступ, правильный ответ --- другой источник, снапшот или обращение к владельцам сервиса за повышенным лимитом.
\end{notebox}

% =====================================================================
\begin{landscape}
\thispagestyle{empty}\pagestyle{empty}
\section{Сводная таблица систем}
\small
\setlength{\tabcolsep}{5pt}
\begin{longtable}{>{\raggedright\arraybackslash}p{2.7cm}>{\raggedright\arraybackslash}p{3.4cm}>{\raggedright\arraybackslash}p{3.0cm}>{\raggedright\arraybackslash}p{6.2cm}>{\raggedright\arraybackslash}p{7.6cm}}
\toprule
\textbf{Система} & \textbf{Область} & \textbf{Ключ} & \textbf{Лимиты (кратко)} & \textbf{Автоматизация и блокировки}\\
\midrule
\endhead
PubMed (E-utilities) & Биомедицина, MEDLINE, MeSH & Необязателен, бесплатный & 3 запр./с без ключа, 10 с ключом; до 10\,000 ID за ESearch & API разрешён. HTML-скрейпинг: блокировка IP («Access Denied», misuse).\\
PubMed Central (PMC) & Полные тексты биомедицинских статей & Не нужен & Через OAI-PMH, FTP, Cloud; E-utilities для метаданных & Систематическая выгрузка с основного сайта запрещена, IP блокируется. Разрешены только OAI-PMH, FTP, Cloud.\\
Europe PMC & Науки о жизни, препринты, патенты & Не нужен & Официальные числа не подтверждены; начинайте с 1--3 запр./с & API открыт; по сообщению стороннего описания, \code{robots.txt} разрешает путь REST API.\\
arXiv & Физика, математика, CS, q-bio, q-fin, stat, econ & Не нужен & 1 запрос/3\,с, одно соединение; не более 30\,000 результатов; срезы до 2000 & API на \code{export.arxiv.org}. Основной сайт: «Robots Beware», при нарушении 403 и блокировка.\\
OpenAlex & Все науки; граф работ, авторов, организаций & \textbf{Обязателен} с 13.02.2026, бесплатный & Дневной бюджет в долларах (с ключом \$1/день; цифры расходятся); до 100 запр./с; одиночные записи бесплатны & API-first, ключ обязателен. Полный снапшот бесплатно (CC0).\\
Crossref & DOI-метаданные, ссылки, лицензии & Не нужен; нужен \code{mailto} & С 21.07.2026: одиночная запись 5/с (public), 10/с (polite); списки 1/с и 3/с & API для машин; 429 при превышении; лимит привязан к e-mail.\\
Semantic Scholar & Все науки, акцент CS/AI; цитаты, TLDR & Бесплатный, рекомендован & С ключом стартово 1 запр./с; без ключа общий пул, часто 429 & Официальный API. Для больших объёмов есть датасеты.\\
CORE & Агрегатор открытого доступа, полные тексты & Бесплатный, желателен & Токены в день, зависят от типа аккаунта (вторичные данные) & API v3; полный текст только с ключом.\\
Unpaywall & Статус OA и ссылки на легальные копии по DOI & Нужен только e-mail & 100\,000 запросов/день (рекомендация) & API разрешён; снапшот для больших объёмов.\\
DOAJ & Рецензируемые OA-журналы и статьи & Не нужен для поиска & 2 запр./с, очередь пачек до 5 & Публичный API, троттлинг.\\
dblp & Информатика, библиография & Не нужен & 429 с \code{Retry-After}; конкретных цифр нет & Скрейпинг разрешён по FAQ, но сторонние клиенты сообщают о bot-check на \code{dblp.org}; есть XML-дамп и SPARQL.\\
NASA ADS / SciX & Астрономия, астрофизика, физика, науки о Земле & Токен бесплатно & 5000 запросов/день & API с токеном; заголовки \code{X-RateLimit-*}.\\
OpenAIRE & Европейский открытый граф научных продуктов & Токен опционален & 60 запр./час без токена, 7200/час с токеном & API разрешён, лицензия CC-BY.\\
Lens.org & Статьи и патенты, связи между ними & Токен, нужен запрос доступа & Зависит от плана (в одном вузовском плане 5000 запросов/мес, 10/мин) & Платный или по заявке; доступ только через токен.\\
BASE & Агрегатор репозиториев (400+ млн документов) & Белый список IP & Не публикуются & Веб-портал открыт, API только для заранее зарегистрированных IP и User-Agent.\\
Scopus & Мульти\-дисци\-плинар\-ный индекс цитирования & Elsevier API key + подписка & Недельные квоты (Search 20\,000/нед, около 9 запр./с) & Только по ключу; полный функционал через IP вуза или institutional token.\\
Web of Science & Мульти\-дисци\-плинар\-ный индекс цитирования & Ключ + платная лицензия & По плану подписки & Starter API входит в подписку Core Collection; Expanded --- платный.\\
IEEE Xplore & Инженерия, электроника, CS & Бесплатный ключ & Около 200 запросов/день, до 200 записей в ответе (вторичные данные) & API по ключу; полный текст только для подписчиков.\\
Springer Nature & Журналы и книги Springer Nature & Бесплатный ключ & Basic: 500 обращений/день, 100/мин (OA API) & Meta и Open Access API; Premium по договору.\\
Google Scholar & Мульти\-дисци\-плинар\-ный, самый широкий охват & Официального API нет & --- & Автоматические запросы запрещены условиями; блокировка и CAPTCHA.\\
eLIBRARY.RU, КиберЛенинка & Русскоязычная литература & Для части функций eLIBRARY нужен apiKey & Не подтверждено & Открытого поискового API не обнаружено; требуется отдельная проверка.\\
\bottomrule
\end{longtable}
\end{landscape}
\pagestyle{fancy}

% =====================================================================
\section{Карточки систем}
% =====================================================================

Для каждой системы указано: на что она ориентирована, как подключиться, какие лимиты действуют, как строится запрос, как она относится к автоматическим клиентам и какова достоверность данных.

% ---------------------------------------------------------------------
\subsection{PubMed и NCBI E-utilities}
% ---------------------------------------------------------------------

\begin{syscard}{PubMed (NLM / NCBI)}
\row{Фокус}{Биомедицина и науки о жизни. Индексирует метаданные (названия, аннотации, MeSH-термины), не полные тексты. Базовый инструмент систематических обзоров в медицине.}
\row{Программный доступ}{Набор E-utilities: \code{esearch} (поиск ID), \code{efetch} (полные записи), \code{esummary} (краткие сведения), \code{elink}, \code{epost}, \code{einfo}. Адрес: \url{https://eutils.ncbi.nlm.nih.gov/entrez/eutils/}. Ключ выдаётся бесплатно в аккаунте NCBI (Settings $\rightarrow$ API Key Management); в запросах стоит указывать параметры \code{tool} и \code{email}.}
\row{Лимиты}{3 запроса/с без ключа, 10/с с ключом. Более высокий лимит для ключа можно запросить письмом в NCBI с обоснованием. При превышении --- HTTP 429. ESearch возвращает не более 10\,000 ID за вызов; для больших выборок используйте history server (\code{usehistory=y}, \code{WebEnv}, \code{query_key}) или дробите запрос по датам.}
\row{Синтаксис запросов}{Булевы операторы \code{AND}, \code{OR}, \code{NOT}; поля в квадратных скобках: \code{[ti]} название, \code{[ab]} аннотация, \code{[tiab]} название или аннотация, \code{[au]} автор, \code{[ta]} журнал, \code{[mh]} термин MeSH (с автоматическим включением узких терминов), \code{[majr]} основной MeSH-термин, \code{[pt]} тип публикации, \code{[dp]} или \code{[pdat]} дата, \code{[la]} язык, \code{[sb]} подмножество (например, \code{free full text[sb]}). Диапазон дат: \code{2020:2026[dp]}. Близость слов: \code{"asthma treatment"[Title:~3]}.}
\row{Подводные камни}{Несуществующий тег поля может быть молча проигнорирован: PubMed включает автоматическое сопоставление терминов (automatic term mapping). Всегда проверяйте поле \code{querytranslation} в ответе ESearch, чтобы убедиться, что запрос понят так, как вы хотели. Очень свежие статьи ещё не получили MeSH, поэтому сочетайте \code{[mh]} с \code{[tiab]}.}
\row{Блокировка веб-агентов}{Для API --- нет, он предназначен для машин. Для веб-интерфейса картина обратная: при подозрении на скрейпинг NCBI блокирует IP-адрес и показывает страницу «Access Denied \ldots{} possible misuse/abuse» (сервер \code{misuse.ncbi.nlm.nih.gov}); в тексте блокировки прямо рекомендуется освоить E-utilities. Вывод: ходить только через E-utilities.}
\row{Статус данных}{Официальные (руководство E-utilities на NCBI Bookshelf, NBK25497 и NBK25499); описания блокировок подтверждены страницами «Access Denied» в поисковой выдаче.}
\end{syscard}

\begin{lstlisting}[caption={PubMed: поиск через ESearch и выгрузка через EFetch с history server}]
import requests, time

BASE = "https://eutils.ncbi.nlm.nih.gov/entrez/eutils/"
COMMON = {"tool": "litsearch", "email": "you@example.org"}   # add "api_key": "..." if you have one

def pubmed_search(term, retmax=200):
    r = requests.get(BASE + "esearch.fcgi", params={
        **COMMON, "db": "pubmed", "term": term, "retmax": retmax,
        "usehistory": "y", "retmode": "json"}, timeout=60)
    r.raise_for_status()
    res = r.json()["esearchresult"]
    print("translation:", res.get("querytranslation"))   # check how PubMed interpreted the query
    return int(res["count"]), res["webenv"], res["querykey"]

def pubmed_fetch(webenv, query_key, start=0, n=200):
    r = requests.get(BASE + "efetch.fcgi", params={
        **COMMON, "db": "pubmed", "WebEnv": webenv, "query_key": query_key,
        "retstart": start, "retmax": n, "rettype": "abstract", "retmode": "xml"}, timeout=120)
    r.raise_for_status()
    return r.text

count, webenv, qk = pubmed_search('"sickle cell disease"[mh] AND CRISPR[tiab] AND 2020:2026[dp]')
for start in range(0, count, 200):
    xml = pubmed_fetch(webenv, qk, start)
    # ... parse xml (e.g. with xml.etree.ElementTree) ...
    time.sleep(0.4)          # stay under 3 requests/second without a key
\end{lstlisting}

% ---------------------------------------------------------------------
\subsection{PubMed Central (PMC)}
% ---------------------------------------------------------------------

\begin{syscard}{PubMed Central (NLM)}
\row{Фокус}{Архив полных текстов биомедицинских и биологических статей; часть материалов в открытом доступе (Open Access Subset и другие наборы с разрешающими лицензиями).}
\row{Программный доступ}{Метаданные можно искать теми же E-utilities (\code{db=pmc}). Для массовой выгрузки текстов разрешены только специальные сервисы: OAI-PMH, FTP и PMC Cloud Service.}
\row{Лимиты}{Действуют общие лимиты E-utilities (см.\ выше). Для FTP и Cloud ограничения определяются условиями набора данных и лицензией конкретных статей.}
\row{Синтаксис запросов}{Как в PubMed (теги полей в квадратных скобках), с несколькими специфичными полями PMC.}
\row{Блокировка веб-агентов}{Да, и документ об авторском праве PMC формулирует это жёстко: систематическая загрузка пакетов статей с основного сайта PMC любым способом запрещена из-за ограничений авторского права; перечисленные три сервиса (OAI-PMH, FTP, Cloud) названы единственными разрешёнными путями автоматической загрузки, и нельзя использовать другие автоматические процессы даже для статей из Open Access Subset. Нарушителей блокируют по IP-диапазону.}
\row{Статус данных}{Официальные (PMC Copyright Notice, процитированный и пересказанный в библиотечных руководствах).}
\end{syscard}

% ---------------------------------------------------------------------
\subsection{Europe PMC}
% ---------------------------------------------------------------------

\begin{syscard}{Europe PMC (EMBL-EBI)}
\row{Фокус}{Науки о жизни: надмножество PubMed, включает PMC, препринты (bioRxiv, medRxiv и др.), патенты, диссертации. Поиск идёт и по метаданным, и по полным текстам открытых статей. Есть Annotations API с текстовой разметкой (гены, заболевания, химические соединения).}
\row{Программный доступ}{REST API: \url{https://www.ebi.ac.uk/europepmc/webservices/rest/}. Ключ не нужен. Методы: \code{search}, поиск по идентификатору, цитирующие и цитируемые работы, ссылки на сторонние ресурсы, полные тексты PMC в XML. Клиенты есть для R (\code{europepmc}) и Python.}
\row{Лимиты}{Подтверждённых официальных числовых лимитов найти не удалось. Разумно начинать с 1--3 запросов/с и увеличивать осторожно, обрабатывая 429 и 5xx с экспоненциальной паузой.}
\row{Синтаксис запросов}{Язык в духе Lucene: \code{TITLE:"gut microbiome"}, \code{DOI:10.1007/bf00197367}, \code{EXT_ID:22246381} (PMID), \code{ISSN:1553-7404}. Поля вроде \code{OPEN_ACCESS}, \code{HAS_FT}, \code{PUB_YEAR}, \code{SRC} описаны в справке Europe PMC; удобно собрать запрос в Advanced Search Query Builder и скопировать строку. По умолчанию включается расширение синонимами (MeSH и UniProt): чтобы совпадать с сайтом, оставляйте его включённым, чтобы получить точный запрос --- отключайте (\code{synonym=FALSE} в R-клиенте).}
\row{Пагинация}{Курсором: параметр \code{cursorMark} (начните с \code{*}), в ответе приходит \code{nextCursorMark}; размер страницы \code{pageSize} до 1000.}
\row{Блокировка веб-агентов}{Для API не обнаружено; по описанию стороннего инструмента, \code{robots.txt} хоста \code{www.ebi.ac.uk} разрешает путь \code{/europepmc/webservices/rest/} и антибот-слоя на нём нет. Прямой проверки я не выполнял.}
\row{Статус данных}{Смешанные: описание функций и синтаксиса из документации клиентов, лимиты не подтверждены.}
\end{syscard}

\begin{lstlisting}[caption={Europe PMC: поиск с курсорной пагинацией}]
import requests

URL = "https://www.ebi.ac.uk/europepmc/webservices/rest/search"
params = {"query": 'TITLE:"long covid" AND OPEN_ACCESS:y',
          "format": "json", "resultType": "core",
          "pageSize": 100, "cursorMark": "*"}
while True:
    d = requests.get(URL, params=params, timeout=60).json()
    for rec in d["resultList"]["result"]:
        print(rec.get("pmid"), rec.get("doi"), rec["title"])
    nxt = d.get("nextCursorMark")
    if not nxt or nxt == params["cursorMark"]:
        break
    params["cursorMark"] = nxt
\end{lstlisting}

% ---------------------------------------------------------------------
\subsection{arXiv}
% ---------------------------------------------------------------------

\begin{syscard}{arXiv (Cornell University)}
\row{Фокус}{Препринты по физике, математике, информатике, количественной биологии и финансам, статистике, электротехнике, экономике. Рецензирования нет; текст и метаданные доступны сразу.}
\row{Программный доступ}{API на хосте \code{export.arxiv.org}: \url{https://export.arxiv.org/api/query}. Ответ в формате Atom XML (JSON нужно получать конвертацией на своей стороне). Для массового сбора предусмотрены OAI-PMH, дамп метаданных и массовый доступ к полным текстам (S3).}
\row{Лимиты}{Условия использования устаревших API (OAI-PMH, RSS и API запросов) требуют не более одного запроса в три секунды и одного соединения; лимит общий для всех ваших машин, наращивание числа машин для обхода запрещено. Руководство пользователя называет потолок 30\,000 результатов на запрос при срезах не более 2000. Новые статьи появляются в API раз в сутки около полуночи, поэтому результаты стоит кэшировать. Метаданные распространяются как CC0, сами файлы статей имеют отдельные права.}
\row{Синтаксис запросов}{Параметр \code{search_query} с префиксами полей: \code{ti:} название, \code{au:} автор, \code{abs:} аннотация, \code{cat:} категория, \code{all:} все поля, \code{submittedDate:[YYYYMMDDHHMM TO YYYYMMDDHHMM]}; операторы \code{AND}, \code{OR}, \code{ANDNOT}; пробелы кодируются как \code{+}, скобки как \texttt{\%28} и \texttt{\%29}, кавычки как \texttt{\%22}. Параметры \code{start}, \code{max_results}, \code{sortBy}, \code{sortOrder}. Пример: \code{au:del_maestro+ANDNOT+ti:checkerboard}.}
\row{Подводные камни}{Ошибка может приходить как запись в Atom-ленте, а не как HTTP-код; успешный разбор XML ещё не значит, что в ленте статьи. Для регулярного мониторинга новых работ лучше использовать RSS по категориям или OAI-PMH, а не опрос поиска.}
\row{Блокировка веб-агентов}{Да, и arXiv формулирует это жёстче всех: страница «Robots Beware» гласит, что неизбирательная автоматическая загрузка с сайта не допускается; сервер ограничен, приоритет отдан людям. Хост мониторится, нарушителям возвращается 403, а продолжение «беглого огня» после отказа трактуется как атака, на которую отвечают без предупреждения. Страница отказа прямо предлагает использовать \code{export.arxiv.org} для автоматического сбора. Если нужен повышенный лимит, нужно писать в поддержку \emph{заранее}.}
\row{Статус данных}{Официальные (Robots Beware, Terms of Use, руководство API). Один сторонний блог упоминает ужесточение политики с 1~октября 2026\,г.; на официальных страницах я этого не подтвердил, поэтому проверьте анонсы arXiv.}
\end{syscard}

\begin{lstlisting}[caption={arXiv: запрос с паузой 3 секунды}]
import requests, time
import xml.etree.ElementTree as ET

NS = {"a": "http://www.w3.org/2005/Atom",
      "os": "http://a9.com/-/spec/opensearch/1.1/"}
HEADERS = {"User-Agent": "litsearch/1.0 (mailto:you@example.org)"}

def arxiv_search(query, start=0, n=50):
    r = requests.get("https://export.arxiv.org/api/query", headers=HEADERS, params={
        "search_query": query, "start": start, "max_results": n,
        "sortBy": "submittedDate", "sortOrder": "descending"}, timeout=60)
    r.raise_for_status()
    root = ET.fromstring(r.text)
    total = int(root.find("os:totalResults", NS).text)
    items = [{"id": e.find("a:id", NS).text,
              "title": " ".join(e.find("a:title", NS).text.split()),
              "published": e.find("a:published", NS).text}
             for e in root.findall("a:entry", NS)]
    return total, items

total, items = arxiv_search("cat:cs.CL AND abs:retrieval")
time.sleep(3.1)     # arXiv terms: one request per 3 seconds, one connection
\end{lstlisting}

% ---------------------------------------------------------------------
\subsection{OpenAlex}
% ---------------------------------------------------------------------

\begin{syscard}{OpenAlex (OurResearch)}
\row{Фокус}{Открытый каталог научных метаданных по всем областям: работы, авторы, источники, организации, темы, финансирование, гранты. Полностью открыт (лицензия CC0). Сегодня один из лучших универсальных источников для первичного поиска и библиометрии.}
\row{Программный доступ}{REST API: \url{https://api.openalex.org/}. Условия изменились в 2026 году: с 13~февраля 2026\,г. для использования API нужен бесплатный ключ (создаётся в аккаунте на openalex.org, передаётся как \code{api_key=}), без ключа допустимы только демонстрационные вызовы. Полный снапшот для скачивания остаётся бесплатным.}
\row{Лимиты}{Введена модель оплаты по использованию: у каждого ключа есть суточный бюджет в долларах. По документации, стоимость вызова зависит от типа: одиночная запись по ID или DOI бесплатна; список с фильтром порядка \$0.0001; поисковый запрос около \$0.001; семантический поиск и полный текст около \$0.01. Общий потолок --- до 100 запросов в секунду независимо от стоимости. Остаток бюджета можно узнать вызовом \code{/rate-limit}. Совет из документации: \code{per_page=100}, чтобы бюджет расходовался экономнее.}
\row{Подводные камни}{Численные значения бесплатного бюджета в разных источниках и версиях документации расходятся (с ключом называют \$1 в день; для безключевого режима встречаются \$0.01, \$0.10 и абстрактные «кредиты»). Руководствуйтесь ответом \code{/rate-limit} для вашего ключа. Многие учебники до 2026 года пишут «ключ не нужен» --- это устарело.}
\row{Синтаксис запросов}{Параметры \code{search} (полнотекстовый поиск по названию, аннотации и др.), \code{filter} (через запятую; например, \code{publication_year:2025,is_oa:true}), \code{select}, \code{group_by}, \code{sort}. Пакетный поиск по DOI: \code{filter=doi:10.1/a|10.1/b} (до 100 значений). Пагинация курсором: \code{cursor=*}, затем \code{meta.next_cursor}. Аннотации приходят в виде инвертированного индекса (\code{abstract_inverted_index}).}
\row{Блокировка веб-агентов}{Для API --- нет: сервис рассчитан на автоматические клиенты, но теперь требует ключ и тарифицирует нагрузку. Для больших объёмов предусмотрен снапшот; парсить сайт openalex.org не нужно.}
\row{Статус данных}{Официальные (документация developers.openalex.org, рассылка пользователей от 13.01.2026), но числа бюджета противоречивы.}
\end{syscard}

\begin{lstlisting}[caption={OpenAlex: фильтрованный поиск с курсором}]
import os, requests

KEY = os.environ["OPENALEX_API_KEY"]
params = {"search": "retrieval augmented generation",
          "filter": "publication_year:2025-2026,is_oa:true",
          "select": "id,doi,title,publication_year,cited_by_count",
          "per_page": 100, "cursor": "*", "api_key": KEY}
n = 0
while params["cursor"] and n < 500:
    d = requests.get("https://api.openalex.org/works", params=params, timeout=60).json()
    for w in d["results"]:
        print(w["doi"], w["publication_year"], w["title"])
        n += 1
    params["cursor"] = d["meta"].get("next_cursor")
\end{lstlisting}

% ---------------------------------------------------------------------
\subsection{Crossref}
% ---------------------------------------------------------------------

\begin{syscard}{Crossref REST API}
\row{Фокус}{Метаданные всех публикаций с DOI, зарегистрированными через Crossref: примерно 180 млн записей. Лучший источник для поиска DOI по библиографической ссылке, для данных о лицензиях, финансировании, ссылках и датах. Аннотации есть далеко не у всех записей, полного текста нет.}
\row{Программный доступ}{\url{https://api.crossref.org/}. Ключ не нужен. Контактный e-mail в параметре \code{mailto=} переводит запрос в «polite»-пул с более высокими лимитами и позволяет Crossref связаться с вами при проблемах. Для платного сервиса Metadata Plus и других интерфейсов (XML API, OAI-PMH) лимиты не менялись.}
\row{Лимиты}{С 21~июля 2026\,г. лимиты зависят от типа запроса. Одиночная запись: 5 запросов/с (public) и 10/с (polite). Списки (запросы с фильтрами и поиском): 1 запрос/с (public) и 3 запроса/с (polite). Лимит теперь считается по адресу e-mail, чтобы нельзя было обойти его множеством IP. Значения приходят в заголовках \code{x-rate-limit-limit}, \code{x-rate-limit-interval} и \code{x-rate-limit-type}. Ранее (с 1~декабря 2025\,г.) вводились также ограничения по числу одновременных запросов.}
\row{Синтаксис запросов}{\code{/works?query=...}, \code{query.bibliographic=...} (поиск по ссылке целиком), \code{query.author=...}, \code{query.title=...}; фильтры \code{filter=from-pub-date:2023-01-01,type:journal-article,has-abstract:true}; \code{rows} до 1000; \code{select=DOI,title}; глубокая пагинация через \code{cursor=*} и \code{next-cursor}.}
\row{Блокировка веб-агентов}{API --- нет, нужно только соблюдать лимиты и указывать e-mail. Превышение приводит к 429 и автоматической блокировке. Ряд сторонних описаний ещё называет «около 50 запросов/с для polite pool»: это устарело и противоречит анонсу 2026 года.}
\row{Статус данных}{Официальные (блог и форум Crossref: анонс от 5.11.2025 и обновление от 21.07.2026).}
\end{syscard}

\begin{lstlisting}[caption={Crossref: поиск по библиографической строке}]
import requests

r = requests.get("https://api.crossref.org/works", params={
    "query.bibliographic": "Attention is all you need Vaswani 2017",
    "rows": 3, "select": "DOI,title,issued,is-referenced-by-count",
    "mailto": "you@example.org"}, timeout=60)
r.raise_for_status()
print(r.headers.get("x-rate-limit-limit"), r.headers.get("x-rate-limit-type"))
for it in r.json()["message"]["items"]:
    print(it["DOI"], it.get("title"))
\end{lstlisting}

% ---------------------------------------------------------------------
\subsection{Semantic Scholar (Academic Graph API)}
% ---------------------------------------------------------------------

\begin{syscard}{Semantic Scholar (Allen Institute for AI)}
\row{Фокус}{Все науки, но особенно силён в информатике и ИИ. Данные о цитированиях и контекстах цитирования, влиятельных цитированиях, краткие TLDR-аннотации, рекомендации похожих статей, ссылки на открытые PDF.}
\row{Программный доступ}{Academic Graph API: \url{https://api.semanticscholar.org/graph/v1}. Ключ передаётся заголовком \code{x-api-key}. Есть Recommendations API и Datasets API (скачиваемые снапшоты). Эндпоинты: \code{/paper/search} (релевантность), \code{/paper/search/bulk} (массовая выгрузка), \code{/paper/batch} (пакетный запрос по ID, до нескольких сотен за вызов), \code{/paper/{id}}, \code{/author/...}.}
\row{Лимиты}{Для ключа стартовый лимит --- 1 запрос/с на всех эндпоинтах. Без ключа действует общий пул для всех безключевых пользователей (в документации указан порядок 1000 запр./с суммарно на всех), и при нагрузке часто приходит 429. Для задач, где этого не хватает, сам сервис советует скачать датасеты и работать локально.}
\row{Подводные камни}{Источники расходятся: в сторонних руководствах встречаются «100 запросов за 5 минут без ключа», «около 20 запросов в минуту», «100 запросов в секунду с ключом»; это либо устарело, либо относится к иным тарифам. Нужно закладывать повтор при 429 и учитывать, что поиск по релевантности ограничен по глубине --- для полной выгрузки используйте bulk-эндпоинт.}
\row{Синтаксис запросов}{\code{query=} (ключевые слова); дополнительные параметры \code{year=2020-2025}, \code{fieldsOfStudy=}, \code{openAccessPdf}, \code{minCitationCount}, \code{venue}; обязательно выбирайте поля параметром \code{fields=title,year,externalIds,abstract,openAccessPdf}, иначе придут только ID и название.}
\row{Блокировка веб-агентов}{Для API --- нет, только троттлинг (429). Рекомендуется всегда использовать ключ, при 429 повторять запрос с задержкой.}
\row{Статус данных}{Часть официальных (страница продукта и руководство), цифры лимитов противоречивые.}
\end{syscard}

\begin{lstlisting}[caption={Semantic Scholar: поиск с ключом и повтором при 429}]
import os, time, requests

H = {"x-api-key": os.environ["S2_API_KEY"]}
def s2_search(q, limit=100, offset=0):
    for attempt in range(5):
        r = requests.get("https://api.semanticscholar.org/graph/v1/paper/search",
            headers=H, params={"query": q, "limit": limit, "offset": offset,
            "fields": "title,year,externalIds,abstract,openAccessPdf,citationCount"},
            timeout=60)
        if r.status_code == 429:
            time.sleep(2 ** attempt); continue
        r.raise_for_status()
        return r.json()
    raise RuntimeError("rate limited")

d = s2_search("dense passage retrieval")
print(d["total"], [p["title"] for p in d["data"][:3]])
time.sleep(1.1)       # 1 request/second with an introductory key
\end{lstlisting}

% ---------------------------------------------------------------------
\subsection{CORE}
% ---------------------------------------------------------------------

\begin{syscard}{CORE (The Open University, Великобритания)}
\row{Фокус}{Крупнейший агрегатор открытого доступа: сотни миллионов исследований из тысяч репозиториев; полные тексты для десятков миллионов статей. Полезен, когда нужен именно доступный текст, а не только метаданные.}
\row{Программный доступ}{API v3: \url{https://api.core.ac.uk/v3}. Регистрация бесплатная на core.ac.uk/services/api; ключ передаётся заголовком \code{Authorization: Bearer ...} или параметром \code{api_key}. Без ключа работают только базовые запросы метаданных, без полного текста. Эндпоинты поиска требуют завершающего слэша: \code{/v3/search/works/}.}
\row{Лимиты}{Лимиты выражаются в «токенах» в день и в минуту. По стороннему описанию: без аутентификации 100 токенов в день и 10 в минуту, личный аккаунт 1000 в день и 25 в минуту, академический 5000 в день; простой запрос стоит 1 токен, скачивание и scroll-пагинация 3--5.}
\row{Синтаксис запросов}{\code{q=} с поддержкой булевых операторов и полей (например, \code{doi:10.1038/nature12373}); параметры \code{limit}, \code{offset}, scroll для глубокой выгрузки. Метод POST принимает JSON с теми же полями.}
\row{Блокировка веб-агентов}{Для API --- нет, используется квота по токенам. Парсить HTML портала не нужно: для массового сбора у CORE есть датасет.}
\row{Статус данных}{Вторичные (описание в навыках для ИИ-агентов и обёртки Apify); числа лимитов следует проверить на странице CORE.}
\end{syscard}

% ---------------------------------------------------------------------
\subsection{Unpaywall}
% ---------------------------------------------------------------------

\begin{syscard}{Unpaywall (OurResearch)}
\row{Фокус}{Не поисковик по темам, а «разрешитель» открытого доступа: по DOI сообщает, есть ли легальная бесплатная версия, где она лежит (репозиторий, препринт, страница издателя) и какого типа OA (gold, green, hybrid, bronze). Незаменим на втором шаге после поиска метаданных.}
\row{Программный доступ}{\code{GET https://api.unpaywall.org/v2/DOI?email=ВАШ_EMAIL}. Ключ не нужен, но e-mail обязателен. Дополнительно есть эндпоинт поиска по названию (\code{/v2/search}, 50 результатов на странице). Для крупных задач --- снапшот базы и Data Feed.}
\row{Лимиты}{Рекомендация: не более 100\,000 запросов в сутки; жёсткого ограничения в секунду не заявлено, но разумно держать 1--10 запросов в секунду. Если нужно больше, сервис советует скачать снапшот.}
\row{Блокировка веб-агентов}{Нет; сервис создан для программного вызова.}
\row{Статус данных}{Официальные (unpaywall.org/api/v2).}
\end{syscard}

\begin{lstlisting}[caption={Unpaywall: найти бесплатную версию по DOI}]
import requests

def free_version(doi, email="you@example.org"):
    r = requests.get(f"https://api.unpaywall.org/v2/{doi}", params={"email": email}, timeout=30)
    if r.status_code == 404:
        return None
    r.raise_for_status()
    d = r.json()
    best = d.get("best_oa_location") or {}
    return {"is_oa": d["is_oa"], "status": d["oa_status"], "pdf": best.get("url_for_pdf")}

print(free_version("10.1038/nature12373"))
\end{lstlisting}

% ---------------------------------------------------------------------
\subsection{DOAJ}
% ---------------------------------------------------------------------

\begin{syscard}{DOAJ (Directory of Open Access Journals)}
\row{Фокус}{Курируемый каталог рецензируемых журналов открытого доступа (порядка 20\,000 журналов) и их статей (миллионы записей). Полезен для проверки «качества» издания и поиска статей только в OA-журналах. Препринты не индексируются.}
\row{Программный доступ}{Публичный API (документация \url{https://doaj.org/api/v4/docs}); поиск без ключа. Эндпоинты вида \code{/api/search/articles/ЗАПРОС} и \code{/api/search/journals/ЗАПРОС}, журнал по ISSN: \code{/api/search/journals/issn:XXXX-XXXX}. Актуальный префикс пути версии сверяйте по документации.}
\row{Лимиты}{Два запроса в секунду на все маршруты; допускаются «всплески», до пяти запросов пользователя ставятся в очередь и обслуживаются по очереди при среднем темпе 2 запроса/с.}
\row{Синтаксис запросов}{База построена на Elasticsearch, поддерживается синтаксис query string (Lucene): \code{title:climate AND author:Smith}; часть возможностей отключена ради производительности. Параметры \code{page}, \code{pageSize}; небольшие размеры страницы (до десятков) работают надёжнее.}
\row{Блокировка веб-агентов}{Нет, только троттлинг по лимиту.}
\row{Статус данных}{Официальные (документация) и вторичные (обёртки).}
\end{syscard}

% ---------------------------------------------------------------------
\subsection{dblp}
% ---------------------------------------------------------------------

\begin{syscard}{dblp computer science bibliography}
\row{Фокус}{Библиография по информатике: конференции, журналы, авторы, с очень чистыми данными об авторах и местах публикации. Полных текстов и аннотаций нет.}
\row{Программный доступ}{API поиска публикаций: \url{https://dblp.org/search/publ/api?q=...&format=json&h=100&f=0} (параметры \code{h} --- число результатов, \code{f} --- смещение); аналогичные для авторов и мест публикации. Также доступны SPARQL-эндпоинт \code{sparql.dblp.org} и полный XML-дамп, синхронный с сайтом.}
\row{Лимиты}{При превышении возвращается пустой документ со статусом 429 и заголовок \code{Retry-After}, указывающий, сколько секунд ждать. Числа не публикуются. Для большого числа запросов FAQ советует скачать дамп.}
\row{Блокировка веб-агентов}{Противоречивая картина. FAQ dblp прямо разрешает краулинг и предлагает простой XML/JSON API; в \code{robots.txt} могут быть ограничения для некоторых разделов. Одновременно авторы клиентских инструментов в 2026 году сообщают, что \code{dblp.org} отдаёт API-клиентам страницу проверки на бота вместо JSON, что вынудило их перейти на SPARQL-эндпоинт; другие отмечают, что dblp «агрессивно» ограничивает частоту и сбрасывает соединения при быстрых последовательных запросах. Рекомендация: держать паузу 3 секунды, обрабатывать 429 и 503, для объёмов использовать дамп или SPARQL.}
\row{Статус данных}{Противоречивые: официальный FAQ против сообщений третьих сторон.}
\end{syscard}

% ---------------------------------------------------------------------
\subsection{NASA ADS / SciX}
% ---------------------------------------------------------------------

\begin{syscard}{NASA Astrophysics Data System (ADS), ныне SciX}
\row{Фокус}{Астрономия, астрофизика, планетология, гелиофизика, смежная физика и науки о Земле. Богатые данные о цитированиях и библиографические коды (bibcode).}
\row{Программный доступ}{\url{https://api.adsabs.harvard.edu/v1/search/query}. Для доступа нужен аккаунт и токен («Generate a new key» в настройках), передаётся заголовком \code{Authorization: Bearer ТОКЕН}. Есть Python-пакет \code{ads}.}
\row{Лимиты}{По умолчанию 5000 запросов в сутки на токен; лимиты заданы отдельно для каждого эндпоинта и сообщаются в заголовках \code{X-RateLimit-Limit}, \code{X-RateLimit-Remaining}, \code{X-RateLimit-Reset}; сброс в полночь UTC. Повышение лимита --- по письму в службу поддержки ADS.}
\row{Синтаксис запросов}{Синтаксис Solr: \code{author:"Hubble, E"}, \code{title:exoplanet*}, \code{abstract:"dark energy"}, \code{year:[2020 TO 2024]}, \code{pubdate:[2023-01-00 TO *]}, \code{"transiting exoplanets" +JWST -Kepler}; параметры \code{fl} (какие поля вернуть), \code{rows}, \code{start}, \code{sort}.}
\row{Блокировка веб-агентов}{Для API нет, но токен обязателен. Не стоит обращаться к интерфейсу сайта автоматически, когда есть API.}
\row{Статус данных}{Официальные (репозиторий adsabs-dev-api) и вторичные (клиенты).}
\end{syscard}

% ---------------------------------------------------------------------
\subsection{OpenAIRE}
% ---------------------------------------------------------------------

\begin{syscard}{OpenAIRE Graph}
\row{Фокус}{Европейский открытый граф научных продуктов: публикации, наборы данных, программное обеспечение, проекты, финансирование, организации. Особенно полезен для связки работ с грантами (EC, UKRI, NIH, Wellcome) и для поиска открытых научных результатов. Граф большой: порядка 350 млн исследовательских продуктов по данным changelog версии 11.}
\row{Программный доступ}{Публичные API (подробно: \url{https://graph.openaire.eu/docs/}) и дампы графа. Для повышенных лимитов нужна авторизация через OpenAIRE AAI: персональный токен доступа (действует около часа, выдаётся по refresh-токену, который живёт месяц) или регистрация сервиса с client credentials.}
\row{Лимиты}{Без аутентификации до 60 запросов в час, с аутентификацией до 7200 запросов в час (скользящее окно один час). Лимиты видны в заголовках \code{x-ratelimit-limit} и \code{x-ratelimit-used}. Ранее действовал переходный период со значением 7200 для всех, теперь для безключевых ограничение жёсткое.}
\row{Лицензия и блокировка}{Данные под CC-BY (нужна ссылка на источник). Это машинный сервис, блокировки HTML не требуются; при превышении --- 429.}
\row{Статус данных}{Официальные (раздел Terms of use документации).}
\end{syscard}

% ---------------------------------------------------------------------
\subsection{Lens.org}
% ---------------------------------------------------------------------

\begin{syscard}{The Lens}
\row{Фокус}{Научные работы и патенты в одной базе со связями между ними (цитирование патентами, поиск предшествующего уровня техники). Полезен для инновационной аналитики и технологического скаутинга.}
\row{Программный доступ}{\url{https://api.lens.org/}. Доступ предоставляется по заявке (Request access), затем в профиле создаётся токен; POST-запросы передают токен как \code{Authorization: Bearer ...}, GET --- параметром. Запросы формулируются как JSON в стиле Elasticsearch (\code{bool}, \code{must}, \code{match}, \code{range}).}
\row{Лимиты}{Зависят от плана. Например, институциональный план одного университета: 5000 запросов в месяц, 500 записей в запросе для научных работ (100 для патентов), 10 запросов в минуту. Остаток сообщается в заголовках \code{x-rate-limit-*}. Для некоммерческих пользователей встречаются «бесплатные» планы, для коммерческих доступ платный, обычно с пробным периодом.}
\row{Блокировка веб-агентов}{Доступ программный только по токену; автоматический обход веб-интерфейса лицензией Lens не предусмотрен.}
\row{Статус данных}{Часть официальных (docs.api.lens.org), числа по планам вторичные и различаются.}
\end{syscard}

% ---------------------------------------------------------------------
\subsection{BASE (Bielefeld Academic Search Engine)}
% ---------------------------------------------------------------------

\begin{syscard}{BASE}
\row{Фокус}{Огромный агрегатор документов открытых репозиториев и цифровых библиотек (порядка 400 млн документов из десятков тысяч источников), собранных по OAI-PMH. Хорош для поиска «серой» литературы и материалов университетских репозиториев.}
\row{Программный доступ}{Есть HTTP-интерфейс (\code{BaseHttpSearchInterface.fcgi}, ответ JSON), но он \textbf{ограничен предварительно зарегистрированными IP-адресами и User-Agent}: нужно заранее запросить доступ у BASE и дать адреса. Веб-портал при этом открыт для людей.}
\row{Лимиты}{Публично не заданы; определяются при согласовании доступа.}
\row{Блокировка веб-агентов}{Да, по IP и User-Agent: запросы вне белого списка отклоняются. Это скорее политика допуска, чем борьба с ботами: выделенный API есть, но по заявке. Без заявки подойдут другие агрегаторы (CORE, OpenAIRE, OpenAlex).}
\row{Статус данных}{Вторичные (описание клиента rbace и каталогов), но с совпадающими утверждениями из нескольких мест.}
\end{syscard}

% ---------------------------------------------------------------------
\subsection{Scopus (Elsevier)}
% ---------------------------------------------------------------------

\begin{syscard}{Scopus Search API}
\row{Фокус}{Коммерческий мультидисциплинарный индекс рефератов и цитирования, широко используемый в наукометрии. Хорош для метрик, профилей авторов, аффилиаций.}
\row{Программный доступ}{\url{https://api.elsevier.com/content/search/scopus}. API-ключ создаётся бесплатно на dev.elsevier.com (заголовок \code{X-ELS-APIKey}), но объём данных зависит от подписки. Elsevier полагается в первую очередь на IP-адреса вуза, прокси не поддерживаются; для удалённой работы нужен «Institutional Token» (по запросу в поддержку).}
\row{Лимиты}{Недельные квоты на ключ, сбрасываются через неделю после первого использования. Для Scopus Search --- 20\,000 запросов в неделю (около 9 запросов/с по стороннему описанию), Abstract Retrieval порядка 10\,000 (в других описаниях 5000), ScienceDirect Search v2 --- 20\,000 в неделю и 2 запроса/с. Один поиск возвращает не более первых 5000 записей. Остаток в заголовках \code{X-RateLimit-*}; при исчерпании --- 429 «QUOTA EXCEEDED».}
\row{Синтаксис запросов}{Поля в стиле \code{TITLE-ABS-KEY("metabolic syndrome") AND PUBYEAR > 2020 AND AFFIL(Oxford)}; также \code{AUTH()}, \code{SRCTITLE()}, \code{DOI()}, \code{LANGUAGE()}. Страницы: \code{count} (до 25 для стандартного представления, до 200 для некоторых), \code{start}.}
\row{Блокировка веб-агентов}{Для API нет блокировки, но жёсткая зависимость от ключа, квот и подписки. Скрейпинг сайта scopus.com лицензией, как правило, запрещён.}
\row{Статус данных}{Официальные (dev.elsevier.com: api\_key\_settings) и вторичные (клиентские библиотеки).}
\end{syscard}

% ---------------------------------------------------------------------
\subsection{Web of Science (Clarivate)}
% ---------------------------------------------------------------------

\begin{syscard}{Web of Science APIs}
\row{Фокус}{Коммерческий индекс цитирования (Core Collection и др.) для наукометрии, оценки журналов и организаций.}
\row{Программный доступ}{Через Clarivate Developer Portal: регистрируете приложение и подписываетесь на нужный API. Два основных продукта: \textbf{Starter API} (поиск и импорт публикаций; входит в подписку на Core Collection) и \textbf{Expanded API} (полные метаданные, адреса, цитирования, финансирование; требует отдельной платной лицензии). У организаций бывают уровни Basic, Intermediate, Premium.}
\row{Лимиты}{Определяются тарифом подписки (число документов и запросов); публичных универсальных чисел нет.}
\row{Синтаксис запросов}{Теги поля вида \code{TS=} (тема), \code{AU=} (автор), \code{PY=} (год), \code{SO=} (источник), \code{DO=} (DOI), с операторами \code{AND}, \code{OR}, \code{NOT}.}
\row{Блокировка веб-агентов}{Доступ программный только по ключу и подписке; обход сайта без договора не предполагается.}
\row{Статус данных}{Официальные (Developer Portal) и вторичные (документация партнёров).}
\end{syscard}

% ---------------------------------------------------------------------
\subsection{IEEE Xplore}
% ---------------------------------------------------------------------

\begin{syscard}{IEEE Xplore Gateway API}
\row{Фокус}{Инженерия, электроника, телекоммуникации, информатика: журналы и труды конференций IEEE и IET, стандарты.}
\row{Программный доступ}{\url{https://ieeexploreapi.ieee.org/api/v1/search/articles}. Бесплатный ключ выдаётся на developer.ieee.org, передаётся параметром \code{apikey}. Метаданные и ссылки доступны всем; полный текст --- для открытых статей или при наличии токена подписки.}
\row{Лимиты}{По стороннему описанию бесплатный уровень допускает около 200 вызовов в день и до 200 записей на вызов (параметр \code{max_records}); при превышении --- 429. Точные условия уточняйте на портале.}
\row{Синтаксис запросов}{Параметры \code{querytext}, \code{author}, \code{article_title}, \code{publication_year}, \code{start_record}, \code{max_records}, \code{sort_field}; булевы операторы в \code{querytext}.}
\row{Блокировка веб-агентов}{API --- по ключу; HTML-страницы издателя с автоматическим доступом не рассчитаны, а лицензионные условия обычно запрещают массовую выгрузку.}
\row{Статус данных}{Вторичные.}
\end{syscard}

% ---------------------------------------------------------------------
\subsection{Springer Nature (Meta API и Open Access API)}
% ---------------------------------------------------------------------

\begin{syscard}{Springer Nature API}
\row{Фокус}{Журналы, книги и главы издательства Springer Nature, включая Nature Portfolio и BMC. Meta API охватывает метаданные всего контента (подписного и открытого), Open Access API --- открытые статьи с полным текстом; для текстового и дата-майнинга предусмотрен отдельный Full Text / TDM API.}
\row{Программный доступ}{Регистрация на dev.springernature.com, ключ нужного типа (Meta, Open Access, TDM). Адрес Meta API: \url{https://api.springernature.com/meta/v2/json}; параметры \code{q}, \code{p} (размер страницы), \code{s} (стартовая позиция), \code{api_key}.}
\row{Лимиты}{По официальной странице: уровень Basic бесплатный, 500 обращений в сутки для Open Access API, троттлинг 100 обращений в минуту, один ключ на аккаунт; уровень Premium (университеты, организации) --- 10\,000 обращений в сутки, 300 в минуту, несколько ключей. Таблица «Managing Rate Limits» даёт сопоставимые числа для Meta API.}
\row{Синтаксис запросов}{Например, \code{q=keyword:"cancer"}, поля \code{title:}, \code{doi:}, \code{journal:}, \code{subject:}, \code{year:}, булевы операторы; результаты --- JSON или XML.}
\row{Блокировка веб-агентов}{API --- нет; для сайта link.springer.com автоматическая массовая загрузка запрещена условиями использования, на практике может блокироваться системами защиты. Используйте API.}
\row{Статус данных}{Официальные (datasolutions.springernature.com, dev.springernature.com).}
\end{syscard}

% ---------------------------------------------------------------------
\subsection{Google Scholar}
% ---------------------------------------------------------------------

\begin{syscard}{Google Scholar}
\row{Фокус}{Самый широкий охват: статьи, препринты, диссертации, книги, патенты, судебные документы; удобные цитирования и метрики авторов и журналов. Но это поисковик для людей, а не источник данных.}
\row{Программный доступ}{\textbf{Официального API нет.} Всё, что называют «Google Scholar API», --- это либо сторонние платные сервисы-скрейперы, либо библиотеки, которые разбирают HTML страницы результатов.}
\row{Условия}{Общие условия Google запрещают доступ к сервисам автоматизированными средствами (роботы, пауки, скрейперы) без явного предварительного разрешения и отправку автоматизированных запросов; для Scholar такое разрешение на массовые запросы не выдавалось.}
\row{Блокировка веб-агентов}{Да, самая известная и жёсткая: после серии запросов с одного адреса появляется страница «unusual traffic from your computer network», затем CAPTCHA, а при продолжении --- блокировка IP на десятки минут и дольше (подтверждено многолетними жалобами пользователей Zotero и других инструментов). Сторонние скрейпинг-сервисы строят бизнес именно на обходе этих защит, что само по себе нарушает условия.}
\row{Рекомендация}{Не использовать для автоматического поиска. Замена по функциям: OpenAlex и Semantic Scholar (широкий охват, цитирования), Crossref (DOI), CORE и Unpaywall (доступ к тексту). Если Scholar нужен, пользуйтесь им вручную, а Zotero и подобным инструментам задавайте паузы и поддерживайте импорт по DOI, где Scholar не требуется.}
\row{Статус данных}{Официальные условия и многократные независимые отчёты пользователей.}
\end{syscard}

% ---------------------------------------------------------------------
\subsection{eLIBRARY.RU и КиберЛенинка}
% ---------------------------------------------------------------------

\begin{syscard}{Российские ресурсы}
\row{eLIBRARY.RU}{Крупнейшая электронная библиотека научной периодики на русском языке; платформа РИНЦ (Russian Science Citation Index) и Science Index. Каждому документу присваивается код EDN, по которому можно ссылаться. В документах платформы упоминается API, открытый для организаций и частных лиц с регистрацией и получением кода доступа \code{apiKey}; из описания следует, что он ориентирован на связь ссылок и цитируемых документов по EDN. Открытого поискового API для сбора библиографии по темам я не нашёл.}
\row{КиберЛенинка}{Бесплатная научная библиотека с полным текстом статей, каталог по классификации OECD; читать можно без регистрации. Документированного публичного API у меня подтвердить не удалось.}
\row{Блокировка веб-агентов}{Данных о политике ботов найти не удалось; при отсутствии документированного API автоматизация HTML-страниц сопряжена с юридическим и техническим риском. Сначала проверьте условия использования и \code{robots.txt} (скрипт в разделе~\ref{sec:probe}), при необходимости напишите владельцам.}
\row{Альтернативы}{Русскоязычные журналы с DOI находятся через Crossref и OpenAlex; репозитории --- через CORE, OpenAIRE и BASE.}
\row{Статус данных}{Слабый: подтверждены только описания продуктов, но не программный доступ.}
\end{syscard}

% =====================================================================
\section{Блокируют ли системы веб-агентов}
% =====================================================================

\subsection{Общая картина}

Короткий ответ: \textbf{сервисы блокируют не «ботов вообще», а автоматические клиенты, которые обходят официальный интерфейс.} Почти везде картина двухслойная:

\begin{enumerate}
  \item \textbf{Официальный API} построен для машин. Его не блокируют, но он держится на лимитах, ключах, квотах и обязательной идентификации (e-mail, \code{tool}, User-Agent с контактом). Нарушение лимита заканчивается 429, а затем временной блокировкой ключа, e-mail или IP.
  \item \textbf{HTML-интерфейс} рассчитан на людей. Здесь блокировки реальны и заметны: PubMed/PMC (IP-блокировка «possible misuse»), arXiv (403 и автоматическая реакция), Google Scholar (CAPTCHA и бан IP), издательские сайты за WAF-системами.
  \item \textbf{Платные индексы} (Scopus, Web of Science, IEEE, Springer) вообще не предполагают доступа без ключа и договора; скрейпинг их страниц нарушает лицензию.
\end{enumerate}

\subsection{Сводка по системам}

\begin{longtable}{>{\raggedright\arraybackslash}p{2.8cm}>{\raggedright\arraybackslash}p{4.3cm}>{\raggedright\arraybackslash}p{5.2cm}>{\raggedright\arraybackslash}p{3.0cm}}
\toprule
\textbf{Система} & \textbf{Машинный доступ} & \textbf{Что происходит при автоматизации HTML} & \textbf{Идентификация}\\
\midrule
\endhead
PubMed & E-utilities & Блокировка IP «misuse/abuse» & \code{tool}, \code{email}\\
PMC & OAI-PMH, FTP, Cloud & Массовая загрузка запрещена, IP-блок & \code{tool}, \code{email}\\
Europe PMC & REST API & Не обнаружено (не проверял) & User-Agent\\
arXiv & \code{export.arxiv.org}, OAI-PMH, S3 & 403 и жёсткая автоматическая реакция на нарушителей & User-Agent\\
OpenAlex & API с ключом, снапшот & Не предполагается; ключ обязателен & \code{api_key}\\
Crossref & REST API & Не предполагается; 429 и блок по e-mail & \code{mailto}\\
Semantic Scholar & Graph API, датасеты & Не предполагается; 429 на общем пуле & \code{x-api-key}\\
CORE & API v3, датасет & Не предполагается; квота токенов & Bearer-ключ\\
Unpaywall & REST по DOI, снапшот & Не предполагается & \code{email}\\
DOAJ & Публичный API & Не предполагается; 2 запр./с & User-Agent\\
dblp & API, SPARQL, XML-дамп & Сообщают о bot-check на API-хосте & User-Agent\\
ADS/SciX & API с токеном & Не предполагается & Bearer-токен\\
OpenAIRE & API, дампы & Не предполагается; 60 запр./ч без токена & Bearer-токен\\
Lens & API с токеном & Только по токену & Bearer-токен\\
BASE & API по белому списку IP & Вне списка отклоняется & IP и User-Agent\\
Scopus, WoS, IEEE, Springer & API по ключу и подписке & Лицензионный запрет; может блокироваться WAF & Ключ\\
Google Scholar & Нет & CAPTCHA, бан IP, нарушение условий & ---\\
eLIBRARY.RU, КиберЛенинка & Не подтверждён & Не изучено & ---\\
\bottomrule
\end{longtable}

\subsection{Что стоит за блокировками на практике}

\begin{itemize}
  \item \textbf{Файл \code{robots.txt}} --- добровольный стандарт, но им пользуются как сигналом. Исследование 2025 года показало, что 71 из 200 самых посещаемых доменов полностью закрыли доступ ботам, а число запретов для конкретных агентов растёт; одновременно часть ботов, особенно ИИ-краулеры, вообще не проверяют \code{robots.txt}, поэтому сайты всё чаще полагаются на сетевую защиту.
  \item \textbf{Сетевая защита уровня Cloudflare и аналогов.} С 2025 года Cloudflare блокирует ИИ-краулеры по умолчанию и поддерживает режим оплаты за обход (ответ HTTP 402); неавторизованные боты получают 403 или проверку. Сообщалось и о побочных эффектах: в августе 2026 года отмечали, что режим блокировки ИИ может давать 403 даже поисковым роботам. Многие издательские сайты работают за подобными фильтрами, поэтому ваш скрипт на облачном сервере с «датацентровым» IP рискует получить отказ.
  \item \textbf{Наблюдаемые отказы.} Когда страницы NCBI/PMC и издательские сайты (например, Wiley) запрашивались автоматическим загрузчиком в поисковой выдаче при подготовке этого документа, многие возвращали «Access Denied». Это косвенное свидетельство, но оно согласуется с описанными политиками.
  \item \textbf{Блокировка бывает ложной.} Часть блокировок срабатывает на общий прокси или NAT организации («иногда блокировку вызывает другой пользователь за тем же прокси», как пишет arXiv), поэтому при работе с университетской сети нужна особая аккуратность.
\end{itemize}

\subsection{Правила, которые снижают риск блокировки}

\begin{enumerate}
  \item Всегда используйте официальный API или дамп; HTML --- в последнюю очередь.
  \item Представляйтесь: User-Agent вида \code{MyTool/1.0 (+https://site; mailto:me@org)}, плюс \code{mailto}, \code{email}, \code{tool}, ключ --- там, где они предусмотрены.
  \item Держите частоту ниже заявленных лимитов, один поток на хост по умолчанию; для arXiv --- не чаще одного запроса в 3 секунды.
  \item Обрабатывайте 429 и 503 с учётом заголовка \code{Retry-After} и экспоненциальной паузой; 403 не повторять --- это сигнал остановиться и разобраться.
  \item Кэшируйте ответы на диск; повторный запрос того же самого не должен идти по сети.
  \item Для объёмов от десятков тысяч записей берите снапшот (OpenAlex, Crossref, Unpaywall, Semantic Scholar, dblp, arXiv).
  \item При необходимости более высокого лимита пишите владельцам сервиса \emph{заранее}: NCBI, arXiv, ADS и другие прямо предлагают такой путь.
  \item Не обходите защиту: не решайте CAPTCHA автоматически, не меняйте IP и подпись клиента, чтобы скрыть автоматизацию.
\end{enumerate}

% =====================================================================
\section{Как выбрать систему под задачу}
% =====================================================================

\begin{longtable}{>{\raggedright\arraybackslash}p{4.4cm}>{\raggedright\arraybackslash}p{5.2cm}>{\raggedright\arraybackslash}p{5.2cm}}
\toprule
\textbf{Задача} & \textbf{Основной источник} & \textbf{Дополнение}\\
\midrule
\endhead
Медицина и биология, систематический обзор & PubMed (E-utilities) + Europe PMC & Unpaywall, PMC OA (через разрешённые сервисы)\\
Универсальный поиск, библиометрия & OpenAlex & Crossref, Semantic Scholar\\
Информатика и ИИ & Semantic Scholar, dblp & arXiv, OpenAlex\\
Физика, математика, свежие препринты & arXiv & ADS (астрофизика), Semantic Scholar\\
Найти DOI по ссылке & Crossref (\code{query.bibliographic}) & OpenAlex\\
Получить бесплатный полный текст & Unpaywall, CORE & Europe PMC, arXiv\\
Цитирования и метрики с подпиской & Scopus, Web of Science & OpenAlex как бесплатная замена\\
Патенты и связь «наука $\leftrightarrow$ патенты» & Lens.org & ---\\
Гранты и финансирование (ЕС) & OpenAIRE & OpenAlex, Crossref\\
Серая литература и репозитории & CORE, OpenAIRE & BASE (по заявке)\\
Журналы открытого доступа и их качество & DOAJ & ---\\
Русскоязычные публикации & Crossref, OpenAlex, CORE & eLIBRARY.RU и КиберЛенинка (проверить условия)\\
\bottomrule
\end{longtable}

\subsection{Рекомендуемая двухэтапная схема}

\begin{enumerate}
  \item \textbf{Поиск кандидатов} (метаданные): запросы в 2--3 системы, подходящие предметной области (например, PubMed + OpenAlex + Semantic Scholar).
  \item \textbf{Нормализация и дедупликация}: приведите DOI к нижнему регистру, убирайте префикс \code{https://doi.org/}; без DOI сопоставляйте по нормализованному названию, году и первому автору. Храните PMID, arXiv-ID, OpenAlex ID как дополнительные ключи.
  \item \textbf{Обогащение}: Crossref (лицензия, финансирование), Semantic Scholar или OpenAlex (цитирования), Unpaywall (OA-ссылка).
  \item \textbf{Получение полного текста} только из разрешённых мест: Europe PMC, PMC (OAI-PMH/FTP/Cloud), arXiv через \code{export.arxiv.org}, ссылки Unpaywall и CORE.
  \item \textbf{Воспроизводимость}: сохраняйте точную строку запроса, дату, версию API, число найденных записей, а для PubMed --- \code{querytranslation}.
\end{enumerate}

% =====================================================================
\section{Универсальный клиент на Python}
% =====================================================================

Клиент ниже централизует три вещи, без которых автоматический поиск быстро упирается в блокировки: паузы по хостам, повтор при 429/503 с уважением к \code{Retry-After} и остановку при 403.

\begin{lstlisting}[caption={Вежливый HTTP-клиент с паузами по хостам, повторами и кэшем}]
import hashlib, json, pathlib, threading, time
from collections import defaultdict
from urllib.parse import urlparse
import requests

USER_AGENT = "LitSearch/1.0 (+https://example.org/litsearch; mailto:you@example.org)"

# Minimum interval between requests, per host (seconds)
MIN_INTERVAL = {
    "eutils.ncbi.nlm.nih.gov": 0.34,      # 3 rps without a key (0.11 with a key)
    "export.arxiv.org": 3.1,              # arXiv: 1 request / 3 s, single connection
    "api.crossref.org": 0.4,              # polite pool, list requests: 3 rps
    "api.openalex.org": 0.2,
    "api.semanticscholar.org": 1.1,       # 1 rps with an introductory key
    "www.ebi.ac.uk": 0.35,
    "doaj.org": 0.55,                     # 2 rps
    "api.unpaywall.org": 0.1,
}
CACHE = pathlib.Path(".http_cache"); CACHE.mkdir(exist_ok=True)
_last = defaultdict(float); _lock = threading.Lock()
_session = requests.Session(); _session.headers["User-Agent"] = USER_AGENT

class Blocked(RuntimeError):
    """403/401 from the server: stop and investigate, do not retry blindly."""

def polite_get(url, params=None, headers=None, retries=5, timeout=60, use_cache=True):
    key = hashlib.sha256(json.dumps([url, params], sort_keys=True, default=str).encode()).hexdigest()
    f = CACHE / key
    if use_cache and f.exists():
        return f.read_text(encoding="utf-8")
    host = urlparse(url).netloc
    for attempt in range(retries):
        with _lock:
            wait = MIN_INTERVAL.get(host, 1.0) - (time.monotonic() - _last[host])
            if wait > 0:
                time.sleep(wait)
            _last[host] = time.monotonic()
        r = _session.get(url, params=params, headers=headers, timeout=timeout)
        if r.status_code in (429, 503):
            try:
                delay = float(r.headers.get("Retry-After", ""))
            except ValueError:
                delay = 2 ** attempt * 2
            time.sleep(delay)
            continue
        if r.status_code in (401, 403):
            raise Blocked(f"{r.status_code} from {host}: check key, terms and robots.txt")
        r.raise_for_status()
        if use_cache:
            f.write_text(r.text, encoding="utf-8")
        return r.text
    raise RuntimeError(f"gave up on {url}")

def norm_doi(d):
    if not d:
        return None
    d = d.strip().lower()
    for p in ("https://doi.org/", "http://doi.org/", "doi:"):
        if d.startswith(p):
            d = d[len(p):]
    return d
\end{lstlisting}

% =====================================================================
\section{Проверка доступности и \texttt{robots.txt} из вашей сети}\label{sec:probe}
% =====================================================================

Так как политики и лимиты меняются, а ответы зависят от вашего IP (университет, облако, домашний провайдер), разумно один раз проверить их у себя. Скрипт ниже: читает \code{robots.txt} выбранных хостов, проверяет, разрешает ли он вашему User-Agent типичный путь, делает по одному лёгкому запросу к API и сообщает код ответа, заголовки лимитов и признаки CAPTCHA или проверки на бота. Он делает минимум запросов и не пытается ничего обойти.

\begin{lstlisting}[caption={probe\_access.py: проверка robots.txt и один безопасный запрос к каждому API}]
import re, time, urllib.robotparser as rp
import requests

UA = "LitSearchProbe/1.0 (mailto:you@example.org)"
TARGETS = [  # (name, robots host, sample path to test in robots, API test URL)
    ("PubMed/E-utilities", "https://eutils.ncbi.nlm.nih.gov", "/entrez/eutils/esearch.fcgi",
     "https://eutils.ncbi.nlm.nih.gov/entrez/eutils/esearch.fcgi?db=pubmed&term=asthma&retmax=1&retmode=json"),
    ("PubMed website", "https://pubmed.ncbi.nlm.nih.gov", "/?term=asthma", None),
    ("arXiv API", "https://export.arxiv.org", "/api/query",
     "https://export.arxiv.org/api/query?search_query=all:test&max_results=1"),
    ("arXiv website", "https://arxiv.org", "/list/cs.CL/recent", None),
    ("Europe PMC", "https://www.ebi.ac.uk", "/europepmc/webservices/rest/search",
     "https://www.ebi.ac.uk/europepmc/webservices/rest/search?query=asthma&format=json&pageSize=1"),
    ("Crossref", "https://api.crossref.org", "/works",
     "https://api.crossref.org/works?rows=1&mailto=you@example.org"),
    ("Semantic Scholar", "https://api.semanticscholar.org", "/graph/v1/paper/search",
     "https://api.semanticscholar.org/graph/v1/paper/search?query=test&limit=1"),
    ("DOAJ", "https://doaj.org", "/api/search/articles/test",
     "https://doaj.org/api/search/articles/test?pageSize=1"),
    ("dblp", "https://dblp.org", "/search/publ/api",
     "https://dblp.org/search/publ/api?q=test&format=json&h=1"),
    ("Google Scholar (robots only)", "https://scholar.google.com", "/scholar?q=test", None),
]
BOT_MARKERS = re.compile(r"captcha|unusual traffic|access denied|verify you are human|robot", re.I)

for name, host, path, api in TARGETS:
    line = {"name": name}
    try:
        p = rp.RobotFileParser(host + "/robots.txt")
        p.read()
        line["robots_allows"] = p.can_fetch(UA, host + path)
        line["crawl_delay"] = p.crawl_delay(UA)
    except Exception as e:
        line["robots_error"] = str(e)[:80]
    if api:
        try:
            r = requests.get(api, headers={"User-Agent": UA}, timeout=30)
            line["status"] = r.status_code
            line["rate_headers"] = {k: v for k, v in r.headers.items()
                                    if k.lower().startswith(("x-rate", "ratelimit", "retry-after"))}
            line["bot_page"] = bool(BOT_MARKERS.search(r.text[:3000])) and r.status_code != 200
        except Exception as e:
            line["request_error"] = str(e)[:80]
    print(line)
    time.sleep(3.2)     # be polite, also to arXiv
\end{lstlisting}

Как читать результат: \code{status 200} и \code{robots_allows True} означают, что для этого пути проблем нет; 403 или \code{bot_page True} --- вас приняли за бота (смените канал: API, снапшот или другая сеть, но не обход защиты); 429 --- снизьте частоту; 401 --- нужен ключ.

% =====================================================================
\section{Чек-лист перед запуском массового поиска}
% =====================================================================

\begin{itemize}
  \item Выбран официальный канал (API, OAI-PMH или снапшот) для каждой системы.
  \item Получены ключи: NCBI (по желанию), OpenAlex (обязательно), Semantic Scholar, CORE, ADS, Scopus и др., где нужны.
  \item В запросах есть идентификация: \code{mailto}, \code{email}, \code{tool}, содержательный User-Agent.
  \item Паузы настроены по самому строгому лимиту хоста (arXiv: 3 с).
  \item Обработаны 429, 503 и \code{Retry-After}; 403 не повторяется.
  \item Включён дисковый кэш; результаты сохраняются сырыми и пересчитываются локально.
  \item Для объёмов от десятков тысяч записей используются снапшоты.
  \item Записаны строка запроса, дата и версия API; для PubMed --- \code{querytranslation}.
  \item Проверено, что полные тексты берутся только из разрешённых источников и с учётом лицензии.
  \item Пройдена проверка \code{probe_access.py} из той сети, где будет работать конвейер.
\end{itemize}

% =====================================================================
\section{Источники}
% =====================================================================

Ниже перечислены страницы, на которых основаны сведения документа (проверены поиском 10~ноября 2026\,г.). Рекомендуется сверять числа с ними перед запуском.

\subsection*{Официальные страницы и документация}
\begin{itemize}[leftmargin=1.2em]
  \item NCBI, Entrez Programming Utilities Help: \url{https://www.ncbi.nlm.nih.gov/books/NBK25497/}, \url{https://www.ncbi.nlm.nih.gov/books/NBK25499/}
  \item NLM Support, API keys и rate limits: \url{https://support.nlm.nih.gov/knowledgebase/article/KA-05318/en-us}
  \item PMC, правила автоматической загрузки: \url{https://www.ncbi.nlm.nih.gov/pmc/about/copyright/}, \url{https://www.ncbi.nlm.nih.gov/pmc/tools/textmining/}
  \item arXiv, Robots Beware: \url{https://info.arxiv.org/help/robots.html}
  \item arXiv, API user manual и bulk data: \url{https://info.arxiv.org/help/api/user-manual.html}, \url{https://info.arxiv.org/help/bulk_data.html}
  \item OpenAlex, документация: \url{https://developers.openalex.org/api-reference/authentication}, \url{https://docs.openalex.org/how-to-use-the-api/rate-limits-and-authentication}
  \item Crossref, анонсы лимитов: \url{https://crossref.org/blog/announcing-changes-to-rest-api-rate-limits/}, \url{https://community.crossref.org/t/refining-rest-api-limits-for-improved-stability-and-reliability/16137}
  \item Semantic Scholar API: \url{https://semanticscholar.org/product/api}, \url{https://webflow.semanticscholar.org/product/api/tutorial}
  \item Unpaywall API: \url{https://unpaywall.org/api/v2}
  \item DOAJ API: \url{https://doaj.org/api/v4/docs}
  \item dblp FAQ о краулинге и лимитах: \url{https://dblp1.uni-trier.de/faq/1474706.html}
  \item NASA ADS, документация API: \url{https://github.com/adsabs/adsabs-dev-api}
  \item OpenAIRE Graph, условия и лимиты: \url{https://graph.openaire.eu/docs/next/apis/terms}
  \item Lens API: \url{https://docs.api.lens.org/getting-started.html}
  \item Elsevier Developer Portal, квоты: \url{https://dev.elsevier.com/api_key_settings.html}
  \item Clarivate Developer Portal: \url{https://developer.clarivate.com/apis/wos}
  \item Springer Nature API: \url{https://datasolutions.springernature.com/products/open-access}, \url{https://dev.springernature.com/docs/rate-limit-details/rate-limits}
\end{itemize}

\subsection*{Вторичные источники (описания клиентов, обзоры, форумы)}
\begin{itemize}[leftmargin=1.2em]
  \item Документация клиентов \code{europepmc} (rOpenSci) и \code{rbace} (BASE).
  \item Описания Apify-актёров для Europe PMC, CORE и DOAJ, содержащие сведения об API и \code{robots.txt}.
  \item Обсуждения блокировок Google Scholar на форумах Zotero и обзор «unusual traffic» от Google.
  \item Руководства пользователей \code{pybliometrics} и \code{scopusflow} (квоты Scopus).
  \item Релизные заметки \code{doi-mcp} и \code{fastpaper-cli} (поведение dblp и Semantic Scholar).
  \item Публикации о блокировке ИИ-краулеров Cloudflare и исследования \code{robots.txt} на arXiv (2503.06035, 2505.21733).
\end{itemize}

\end{document}
